\chapter{初等数学工具箱}

本附录是一份“出发前检查清单”。正文默认你大致见过下面这些工具，但不要求熟练；如果哪一节读起来陌生，在这里花十分钟补一补就够了。每项工具我们只回答三个问题：它是什么、它解决了什么麻烦、一个典型例子。这份清单大致按工具在正文中出场的频率排列，越靠前的越常用。读法建议是“用时再查”，不必预先通读——但每一项后面的例子值得动手验算一遍，验算通过的，才算真正装进了你的工具箱。

\section{分数与比例：整体的一部分}

分数 $\dfrac{a}{b}$ 表示把单位“$1$”平均分成 $b$ 份，取其中 $a$ 份。它解决的麻烦是：整数不够用了——三个人分两张饼，每人分到 $\dfrac{2}{3}$ 张，这个数用整数写不出来。分数运算的两条主规则是：加减先通分，$\dfrac{1}{2}+\dfrac{1}{3}=\dfrac{3}{6}+\dfrac{2}{6}=\dfrac{5}{6}$；乘除很直接，$\dfrac{2}{3}\times\dfrac{3}{5}=\dfrac{2}{5}$，除以一个分数等于乘以它的倒数。

比例是两个量的相对大小。地图上“$1:50000$”表示图上 $1$ 厘米对应实地 $50000$ 厘米，即 $500$ 米：图上量得 $3.6$ 厘米，实地就是 $3.6\times500=1800$ 米。百分比只是把分母固定成 $100$ 的比例：打八折就是付原价的 $80\%$。正比例 $y=kx$ 意味着一个量翻倍另一个也翻倍；反比例 $y=\dfrac{k}{x}$ 意味着一个翻倍另一个减半——速度翻倍，行同样路程的时间就减半。

用百分比描述变化时，务必看清“以谁为整体”：价格先涨 $10\%$ 再降 $10\%$，并不会回到原价——涨的时候基数小，降的时候基数大，最后只剩原价的 $99\%$。生活中大量的“数字陷阱”都藏在这一步里。

\section{方程：把未知数当已知数来用}

方程的核心思想只有一句话：\textbf{给未知数起个名字，然后把它当已知数一样运算}。“某数的 $3$ 倍加 $5$ 等于 $20$”，设某数为 $x$，写成 $3x+5=20$。解方程靠的是“天平原则”：等式两边同时加、减、乘、除以同一个数（除数不为 $0$），等号仍然成立。于是 $3x=15$，$x=5$。

再看一个完整的解题流程：“父亲今年 $40$ 岁，儿子 $12$ 岁，几年后父亲年龄是儿子的 $2$ 倍？”设 $x$ 年后，列方程 $40+x=2(12+x)$，展开得 $40+x=24+2x$，移项得 $x=16$。回代检验：$16$ 年后父亲 $56$ 岁、儿子 $28$ 岁，恰好两倍。设未知数、列方程、解方程、回代检验，四步缺一不可，而最后一步最容易被偷懒跳过。

二元一次方程组是两个天平一起用。例如
\[\begin{cases}x+y=10,\\ x-y=4,\end{cases}\]
两式相加消去 $y$：$2x=14$，$x=7$，代回得 $y=3$。一元二次方程 $x^2-5x+6=0$ 常常可以因式分解为 $(x-2)(x-3)=0$，于是 $x=2$ 或 $x=3$；分解不开时还有求根公式 $x=\dfrac{-b\pm\sqrt{b^2-4ac}}{2a}$ 兜底。

\section{坐标：给平面上的每个点一个地址}

在平面直角坐标系里，每个点对应一对有序实数 $(x,y)$：先向东（负则向西）走 $x$，再向北（负则向南）走 $y$。这把“图形”翻译成了“数”，从此几何问题可以用代数计算，代数式子也可以画成图来看。

\begin{center}
\begin{tikzpicture}[scale=0.8]
\draw[->] (-1,0) -- (5,0) node[right] {$x$};
\draw[->] (0,-1) -- (0,4) node[above] {$y$};
\foreach \t in {1,2,3,4} {\draw (\t,0.08) -- (\t,-0.08);}
\foreach \t in {1,2,3} {\draw (0.08,\t) -- (-0.08,\t);}
\fill (3,2) circle (2pt) node[above right] {$(3,2)$};
\draw[dashed] (3,2) -- (3,0);
\draw[dashed] (3,2) -- (0,2);
\end{tikzpicture}
\end{center}

两点 $(x_1,y_1)$ 与 $(x_2,y_2)$ 之间的距离公式 $\sqrt{(x_1-x_2)^2+(y_1-y_2)^2}$ 其实就是勾股定理穿上坐标的外衣：横差和纵差是两条直角边。中点公式 $\left(\dfrac{x_1+x_2}{2},\dfrac{y_1+y_2}{2}\right)$ 说的是“两个地址各取平均”：$(1,2)$ 与 $(3,6)$ 的中点是 $(2,4)$。一次函数 $y=2x+1$ 的图像是一条直线，斜率 $2$ 表示“$x$ 每增加 $1$，$y$ 增加 $2$”。

\section{三角形：平面几何的基本积木}

三条边围成的三角形是多边形里最“倔强”的：三边长度定了，形状就定了（这就是为什么房梁和自行车架爱用三角形）。三条最常用的性质：

\begin{itemize}
\item 内角和：三个内角之和恒为 $180^\circ$；
\item 勾股定理：直角三角形两直角边 $a,b$ 与斜边 $c$ 满足 $a^2+b^2=c^2$，例如 $3^2+4^2=5^2$；
\item 三边关系：任意两边之和大于第三边，所以 $2,3,6$ 围不成三角形（$2+3<6$）。
\end{itemize}

三角形面积等于“底乘高的一半”，这条公式将来会在积分思想里以“切成细条再加起来”的面目重新出现。两个三角形“全等”指完全重合（判定如“边边边”），“相似”指形状相同、大小可以不同，对应边成比例。相似还有一条常用推论：对应边之比为 $k$ 时，面积之比是 $k^2$——边长翻倍，面积变四倍，这就是为什么大份披萨往往比小份“划算得不成比例”。测量旗杆高度用的正是相似：同一时刻，旗杆高与影长之比等于你的身高与影长之比。

\section{指数与对数：乘法的加速器与倒车挡}

指数把“连乘”缩写：$2^{10}$ 表示十个 $2$ 相乘，等于 $1024$。指数法则 $a^m\cdot a^n=a^{m+n}$ 说的是“连乘的个数相加”；特别地，$a^0=1$（一个也不乘），$a^{-1}=\dfrac{1}{a}$（倒过来）。科学记数法是指数的日常应用：光速约 $3\times10^{8}$ 米每秒，一个氢原子的质量约 $1.67\times10^{-27}$ 千克——太大和太小的数，都靠指数才写得下。

对数是指数的“倒车挡”：$\log_a b$ 回答的问题是“$a$ 的几次方等于 $b$”。因为 $2^{10}=1024$，所以 $\log_2 1024=10$。它解决的麻烦是：乘法太多算不动。利用 $\log_a(xy)=\log_a x+\log_a y$，对数能把乘法变成加法——在没有计算器的年代，天文学家靠对数表把几个月的乘法压缩成几天的加法。

\begin{shuyu}{对数}
直观解释：一个数“需要乘几次底数才能得到”。正式定义：若 $a^x=b$（$a>0$，$a\neq1$，$b>0$），则 $x=\log_a b$。例子：$\log_{10}1000=3$，因为 $10^3=1000$。
\end{shuyu}

\begin{toolbox}
六件随身工具：分数与比例（部分与相对大小）、方程（未知数当已知数用）、坐标（图形与数互译）、三角形（内角和、勾股定理、三边关系）、指数（连乘的缩写）、对数（指数的反问）。正文任何一处用到它们，都可以回到本附录对照。
\end{toolbox}

\section*{思考题}

\textbf{观察题}
\begin{sikao}
\item 一杯糖水含糖 $10\%$，喝掉一半后再加满水，浓度变成多少？提示：跟踪“糖还剩多少”，而不是跟踪“水加了多少”。
\end{sikao}

\textbf{证明题}
\begin{sikao}
\item 证明：边长为 $1$ 的正方形，对角线长为 $\sqrt{2}$。提示：对角线把正方形切成两个什么样的三角形？
\end{sikao}

\textbf{挑战题}
\begin{sikao}
\item 不解方程，判断 $x^2+x+1=0$ 有没有实数解。提示：考察求根公式里根号下的 $b^2-4ac$。
\end{sikao}
